\section*{Introducción: determinación estructural y recuperación de información}
\addcontentsline{toc}{section}{Introducción: determinación estructural y recuperación de información}
\label{libro:prologo}
La extrema y media razón holográfica plantea una cuestión precisa: cómo puede una estructura producir una evaluación numérica, transformarse y conservar la información necesaria para reconocer e invertir esa transformación. El problema comprende tanto la formación del valor como la persistencia de sus relaciones. Una igualdad entre cantidades permite comprobar un balance; una ley de recuperación determina además qué estado lo realiza, qué datos lo individualizan y mediante qué operación se restituye. El presente trabajo desarrolla esa segunda exigencia desde una construcción aritmética de caminos y la lleva hasta sus representaciones lineales, incidenciales y físicas.

El primer resultado que organiza el argumento es la determinación de la constante holográfica de acoplamiento aritmético--geométrico. Su representación ordenada es el registro dodecafásico \(K\); su evaluación posicional es el racional \(\kappa\). La lectura conserva exactamente el registro cuando se fijan la base, la longitud y el origen de fase. El interés del objeto reside también en sus operaciones: sus desplazamientos cíclicos forman una base, sus respuestas permiten identificar operadores y sus imágenes incidenciales transportan las relaciones de la configuración inicial. Estructura, ecuación y valor aparecen así como momentos de una construcción cuya información recuperable se determina matemáticamente.

El segundo resultado proporciona la ley general de esa conservación. Dos transportes isométricos con dominio y codominio comunes producen dos componentes cuyo balance ponderado restituye exactamente el producto interior inicial. La especialización nonádica conserva una norma ternaria de valor \(73\) y determina una transformación de observables con pesos \(72/73\) y \(1/73\). Una política de prolongación y archivo convierte esa identidad local en recuperación estable a profundidad arbitraria. La tesis del artículo adquiere con ello un contenido verificable: el valor evaluado, el estado conservado y la lectura observada se relacionan por mapas explícitos y fórmulas inversas.

\subsection*{El núcleo generador y sus dos realizaciones}
La construcción se establece en un núcleo formal denominado \emph{Holografía Modular Triádica}, en adelante HMT. Su primer objeto es \emph{All Possible Paths}, en adelante APP: una estructura aritmética de caminos sobre el soporte toroidal \((\mathbb Z/9\mathbb Z)^2\), cuyo grafo de Cayley describe los desplazamientos cardinales. Sobre ese mismo soporte se distinguen la evaluación aditiva y la multiplicativa. La conservación del residuo junto con el cociente mantiene el entero evaluado y permite seguir la procedencia de cada publicación. La suma acumula; el producto organiza la segunda hoja mediante su propia evaluación y sus reglas de transporte.

El segundo objeto, el \emph{TRIT}, aporta la firma orientada \(\{-1,0,+1\}\) y distingue los regímenes locales. La ecuación de Euler extendida expresa, en sus álgebras cuadráticas, los regímenes elíptico, parabólico e hiperbólico mediante una ley exponencial común. El tercer objeto es el \emph{Topological Phase Kernel}, en adelante TPK. Su acción compone selección, transporte y actualización de memoria. Ruta, dirección, hoja, residuo, cociente, acarreo y frontera se conservan en el estado enriquecido pertinente; el retorno de fase prolonga su historia. La conexión y holonomía nonádicas pertenecen a esta dinámica.

De ese estado procede la estructura discreta conjunta del continuo y sus dos realizaciones: la contracción de cilindros compatibles determina coordenadas aritméticas, mientras que la incidencia de frontera conserva relaciones combinatorias y geométricas. Las cinco construcciones correlativas del continuo permanecen en ese mismo desarrollo. Los espacios lineales, las formas y los observables utilizados después son realizaciones de objetos ya generados, con sus dominios y normalizaciones. La continuidad del argumento reside en que cada operación conserva las relaciones que necesita la siguiente.

\subsection*{Completación de la unidad y refinamiento reversible}
La completación de una célula de nueve posibilidades por coordenada añade una posición de frontera distinguida. En tres coordenadas, la distribución según el número de posiciones de frontera da
\[
10^3=729+243+27+1=729+271.
\]
El interior, los estratos intermedios y el ápice componen una misma unidad normalizada. En dimensión finita arbitraria, la expansión de \((9+1)^d\) expresa la misma estratificación; sumar las contribuciones de la coordenada omitida demuestra su compatibilidad con la proyección.

El bloque central \(\bigl(\begin{smallmatrix}7&2\\2&7\end{smallmatrix}\bigr)\) proporciona un antecedente aritmético concreto. Sus autovalores son \(9\) y \(5\), y sus concatenaciones satisfacen \(72+27=77+22=99\). La restitución de una unidad y el refinamiento siguiente forman la cadena \((72,27)\mapsto(73,27)\mapsto(729,271)\). La primera operación incorpora la unidad registrada; la segunda cambia la escala y redistribuye una unidad mediante una corrección de suma cero. Es esta secuencia explícita la que permite relacionar centro, completación y proporción conservando la procedencia de cada término.

El refinamiento con acarreo introduce una evolución de esa distribución. La transformación \((x,y)\mapsto(Bx-c,By+c)\) conserva la suma después de multiplicarla por \(B\). Sobre etiquetas enteras admisibles, con \(0\leq c<B\), la división euclídea de la segunda componente recupera primero \(c\) y después el par precedente. La profundidad fija cuántas veces debe aplicarse esa inversa. Cuando se permiten acarreos levantados, su cociente adicional forma parte del registro conservado. La reversibilidad depende, por tanto, de datos precisos y tiene una operación de recuperación explícita.

Esta distinción permite dar un sentido operativo a la conservación holográfica. La lectura normalizada determina una coordenada; el registro conserva las relaciones que permiten reconstruir el estado o continuar su refinamiento. El desarrollo posterior hace lineal esa correspondencia: la biyección entre etiquetas admisibles induce un operador unitario entre sus bases y permite transportar amplitudes, correlaciones y observables.

\subsection*{La función aritmética e incidencial de \(K\)}
El registro \(K\) se obtiene del estado terminal mediante sus visitas, trazas, recuentos, diferencias y carga. La evaluación
\[
\kappa=\frac{N_K}{1000^{12}-1}
\]
es reversible en la carta fijada: multiplicar por el denominador restituye el entero y su separación en doce bloques recupera el registro completo, incluidos los ceros iniciales. La historia enriquecida tiene, además, las coordenadas de profundidad y transporte que utiliza cada prolongación. La recuperación de cada objeto se formula sobre los datos que lo individualizan.

En la lectura aritmética, \(\kappa\) interviene en
\[
\alpha_A=\pi_{\rm HMT}+e_{\rm HMT}-\varphi_{\rm HMT}-4-\kappa.
\]
Las tres coordenadas regionales y el registro proceden del mismo núcleo generador y actúan según su orden constructivo. La carta analítica completa de \(\alpha\) representa esa misma precoordenada mediante una raíz simple y una familia compatible de intervalos. En la lectura incidencial, el registro transporta sus relaciones de fase a la configuración excepcional. La unidad del objeto se manifiesta en la composición efectiva de sus lectores, sus inversas y sus acciones.

La órbita cíclica de \(K\) añade un resultado de identificación. Sus doce vectores forman una base de la realización lineal de doce componentes. Las respuestas a esa base determinan cualquier operador lineal de la clase considerada; para un operador que conmuta con el desplazamiento, una respuesta vectorial completa determina las restantes. El menor autovalor de Gram, \(589609\), proporciona una cota exacta de estabilidad. Dos observaciones complementarias de memoria recuperan también la parte centrada del registro; la carga uniforme completa su reconstrucción. La constante funciona así como un objeto de acoplamiento con capacidad de evaluación, transporte e identificación cuantificada.

\subsection*{El balance bilateral y la transformación de observables}
El entero \(73\) aparece como norma de dos elementos conjugados y como índice de una inclusión reticular de rango dos. La identidad \(10\cdot73=729+1\) relaciona esa norma ternaria con la completación decimal. El teorema bilateral extiende el balance a dos isometrías \(B,C\) entre espacios con dominio y codominio comunes. Para \(a>0\), las componentes \(aB-C\) y \((a+1)B+C\), ponderadas por \(a+1\) y \(a\), conservan el producto interior tras la normalización indicada. La expansión y cancelación de los términos cruzados prueban la identidad en dimensión arbitraria. La especialización \(a=8\) mantiene los operadores nonádicos y su información aritmética.

La conservación del estado completo permite calcular el efecto de una observación particular. Si el transporte relativo \(R\) es unitario y \(G\) es el observable expresado en la carta común, la lectura diagonal transforma
\[
\mathcal T(G)=\frac{72}{73}G+\frac1{73}R^*GR.
\]
Un observable permanece invariante exactamente cuando conmuta con \(R\). El observable transportado sobre el espacio completo conserva toda la lectura inicial. El resultado precisa, en una misma construcción, la información que persiste y el cambio introducido al observarla a través de un canal determinado.

\subsection*{Recuperación a profundidad arbitraria}
En cada paso se puede prolongar una componente y archivar su complemento. El balance finito suma las normas cuadradas de las memorias y del estado terminal. Una política explícita impone la cota geométrica \((729/1241)^N\) sobre la norma cuadrada del operador terminal nonádico. En el límite, toda la norma queda contenida en el archivo infinito; este archivo es una isometría y su adjunto reconstruye el estado con norma operatoria uno.

Los primeros bloques archivados admiten asimismo una inversa exacta corregida por su operador de Gram, con una cota de estabilidad uniforme para profundidades positivas. El archivo infinito, la inversa finita y la aproximación obtenida al truncar el adjunto tienen fórmulas y estimaciones propias. En las dinámicas generales donde subsiste un terminal, su contribución se conserva junto con la memoria. La política demostrada identifica el caso en que ese terminal tiende a cero.

El levantamiento unitario del refinamiento se compone con el balance bilateral. El ejemplo decimal que lleva \((73,27,1)\) a \((729,271)\) mantiene explícita la procedencia del acarreo; su comparación con un acarreo vecino determina el cambio exacto de una lectura. La conservación evolutiva de la unidad queda expresada mediante una secuencia de transformaciones recuperables.

\begin{figure}[htbp]
\centering
\begin{tikzpicture}[>=Stealth,every node/.style={font=\small,align=center},
block/.style={draw=ink,rounded corners=2pt,inner sep=6pt,text width=35mm},
wide/.style={draw=ink,inner sep=6pt,text width=47mm}]
\node[block] (d) {Estado aritmético\\admisible\\\((x,y,c)\)};
\node[block,right=14mm of d] (e) {Estado refinado\\\((10x-c,10y+c)\)};
\node[block,right=14mm of e] (j) {Bases de etiquetas\\\(J\) unitario};
\draw[<->] (d)--node[above,font=\scriptsize]{acarreo} (e);
\draw[->] (e)--(j);
\node[wide,below=17mm of e] (w) {Transporte bilateral\\\(W^*W=I\)};
\draw[->] (j.south)|-(w.east);
\node[block,below left=17mm and 4mm of w] (a) {Lectura conjunta\\\(\frac{72G+R^*GR}{73}\)};
\node[block,below right=17mm and 4mm of w] (b) {Observable transportado\\\(WGW^*\)};
\draw[->] (w)--(a);\draw[->] (w)--(b);
\node[font=\small,below=8mm of a] {Invariante si \(GR=RG\)};
\node[font=\small,below=8mm of b] {Recuperación mediante \(W^*\)};
\end{tikzpicture}
\caption{Composición de operaciones demostradas. El refinamiento conserva la suma normalizada y permite recuperar la etiqueta. Su levantamiento actúa sobre amplitudes de estados y conserva la norma. La lectura de un observable y su transporte completo conservan funciones diferentes, especificadas por las dos fórmulas inferiores.}
\label{libro:fig:composicion}
\end{figure}


\subsection*{Formas, acción y observación cuántica}
El transporte del producto interior se prolonga a cada grado exterior admisible. Áreas y volúmenes conservan también las contribuciones mixtas de lectura y memoria. La incidencia excepcional transporta estas identidades sobre sus imágenes especificadas; las cargas variacionales se obtienen de una acción declarada y de sus simetrías. Las realizaciones electromagnéticas y gravitatorias mantienen las escalas y operadores que intervienen en sus pruebas.

La realización gaussiana muestra por qué conservar el estado y conocer una observación reducida son cuestiones complementarias. En el protocolo declarado de dos entradas gaussianas puras, centradas e idénticas y un número finito de modos, el transporte simpléctico conserva la estructura global, mientras su reducción adquiere un espectro determinado por la componente antilineal de la memoria. El lazo elíptico--parabólico estudiado produce pureza \(9/11\), espectro \((9/10)10^{-j}\) y entropía adimensional \((10/9)\log10-\log9\), con todas las ocupaciones del espacio de Fock conservadas en el cálculo.

El argumento concluye, por ello, con una ley de transformación más rica que una igualdad escalar: se determina el valor, se conserva la estructura que interviene en su lectura y se calculan las condiciones de su recuperación. Las fórmulas de inversión, las cotas de estabilidad y las realizaciones especificadas expresan conjuntamente el significado matemático de la extrema y media razón holográfica.
